\documentclass[11pt]{article}

\usepackage[T1]{fontenc}
\usepackage{amsmath,amssymb,amsthm}
\usepackage{hyperref}

\title{RMSEA under FIML: Local Asymptotics with Missing Data}
\author{}
\date{\today}

\newcommand{\E}{\mathbb{E}}
\newcommand{\tr}{\operatorname{tr}}
\newcommand{\Var}{\operatorname{Var}}
\newcommand{\Cov}{\operatorname{Cov}}
\newcommand{\vech}{\operatorname{vech}}
\newcommand{\rank}{\operatorname{rank}}
\newcommand{\argmin}{\operatorname*{arg\,min}}
\newcommand{\plim}{\operatorname*{plim}}
\newcommand{\Gmis}{\Gamma_{\mathrm{mis}}}
\newcommand{\dto}{\xrightarrow{d}}
\newcommand{\pto}{\xrightarrow{p}}
\newcommand{\rmsea}{\widehat{\varepsilon}}

\theoremstyle{plain}
\newtheorem{theorem}{Theorem}
\newtheorem{lemma}{Lemma}
\newtheorem{corollary}{Corollary}
\theoremstyle{remark}
\newtheorem*{remark}{Remark}

\begin{document}
\maketitle

\section*{Purpose}

Reported FIML RMSEA reuses the complete-data formula verbatim:
$\rmsea=\sqrt{\max(T-df,0)/(n\,df)}$, with $T$ the FIML likelihood-ratio
statistic. Two things change underneath, and both matter.

\begin{enumerate}
  \item \textbf{The population target moves with the design.} RMSEA estimates
        $\varepsilon=\sqrt{F_0/df}$, but under missing data $F_0$ is read in the
        \emph{observed-information} metric, which averages over missingness
        patterns and is dominated by the complete-data information. The same
        structural misspecification therefore produces a \emph{smaller}
        population RMSEA the more data are missing (Theorem~\ref{thm:atten}).
        FIML RMSEA confounds misfit with information loss; it is not a pure model
        property.
  \item \textbf{The robust correction needs the FIML sandwich.} The
        scaling/mixture that brings $T$ back to a $\chi^2$ reference is built
        from $U\Gmis$ with $\Gmis=V^{-1}JV^{-1}$, where the observed-data
        curvature $V$ and the meat $J$ both carry the pattern structure, and
        where (unlike complete data) the observed-versus-expected information
        choice is live even under normality.
\end{enumerate}

This is the missing-data companion to the complete-data RMSEA note: the
local-alternative scaffolding is identical, with the substitution $V\to$ the
pattern-averaged FIML curvature and $\Gamma\to\Gmis$. It also presupposes the
point-estimate consistency of the FIML companion note (FIML targets $\eta_0$
only under the MAR conditional-moment conditions); if those fail, $F_0$ is not
even the right population number to root.

\section{The FIML statistic and its information objects}

Row $i$ has observed index set $o_i\subseteq\{1,\dots,p\}$. The saturated
(H1) working model is Gaussian on the observed subvector,
\[
  \ell_i(\eta)
    = \log\phi_{|o_i|}\{x_{i,o_i};\,\mu_{o_i},\Sigma_{o_io_i}\},
  \qquad
  \eta=(\mu,\vech\Sigma)\in\mathbb{R}^{p^\ast},\ \ p^\ast=\tfrac{p(p+3)}{2}.
\]
Missing data forces a mean structure, so $\eta$ \emph{always} carries the $p$
means; $df=p^\ast-t$ counts them. Let $s_i=\partial\ell_i/\partial\eta$ be the
casewise score (sparse in $\eta$: supported on the moments pattern $o_i$
observes). Define the two saturated population objects
\[
  V = \plim\,\tfrac1N\textstyle\sum_i\big(-\partial^2\ell_i/\partial\eta\,\partial\eta'\big),
  \qquad
  J = \plim\,\tfrac1N\textstyle\sum_i \Cov(s_i),
\]
the observed-data normal \emph{curvature} (expected information per case,
averaged over the pattern distribution) and the score \emph{meat}. The saturated
FIML estimator obeys $\sqrt N(\hat\eta_{\mathrm{sat}}-\eta_0)\dto
N(0,\Gmis)$ with the sandwich covariance $\Gmis=V^{-1}JV^{-1}$. The structural
model is $\theta\mapsto\sigma(\theta)\equiv\eta(\theta)$ with full-rank Jacobian
$\dot\sigma=\partial\eta/\partial\theta'$, and the residual-information projector
is, exactly as in the complete-data note,
\begin{equation}\label{eq:U}
  U = V - V\dot\sigma(\dot\sigma'V\dot\sigma)^{-1}\dot\sigma'V,
  \qquad U\dot\sigma=0,\ \ \rank U = df .
\end{equation}
The FIML LRT is $T=2\{\ell_{\mathrm{sat}}(\hat\eta_{\mathrm{sat}})
-\ell_{\mathrm{model}}(\hat\theta)\}$; a second-order expansion gives, with
$w=\sqrt N(\hat\eta_{\mathrm{sat}}-\eta(\theta^\ast))$,
\begin{equation}\label{eq:Tquad}
  T = w'Uw + o_p(1).
\end{equation}
Everything below is a statement about the quadratic form \eqref{eq:Tquad} with
$w\dto N(\tau,\Gmis)$.

\section{Local asymptotics: the noncentral chi-square}

The drift is the same device as in the complete-data note: a fixed off-model
direction shrinking at the $\sqrt N$ rate, $\Sigma_N=\Sigma(\theta^\ast)
+\Delta/\sqrt N$ (plus the matching mean drift), with limiting model error
$\tau=\plim w$ fixed.

\begin{theorem}[Noncentral limit under FIML]\label{thm:ncp}
Under the drift, multivariate normality, MCAR, and identification (every
$\eta$-coordinate seen by some pattern), the information identity $J=V$ holds, so
$\Gmis=V^{-1}$ and $U\Gmis=UV^{-1}$ is idempotent of rank $df$. Hence
\[
  T \dto \chi^2_{df}(\lambda),
  \qquad
  \lambda = \tau'U\tau = \lim_N N F_0 .
\]
\end{theorem}

\begin{proof}
Under MCAR + normality the saturated FIML score is a correctly specified
likelihood score, so $J=V$ (information identity) and $\Gmis=V^{-1}$. Then
$UV^{-1}$ is idempotent of rank $df$ by the computation of the complete-data note
(replace its $V$ by the present curvature). A quadratic form $w'Uw$ with
$w\sim N(\tau,V^{-1})$ and $UV^{-1}$ idempotent of rank $df$ is
$\chi^2_{df}(\lambda)$, $\lambda=\tau'U\tau$. The population discrepancy expands
in the same $V$-metric, giving $NF_0\to\tau'U\tau$.
\end{proof}

So the noncentrality interpretation survives: $\rmsea=\sqrt{\max(T-df,0)/(n\,df)}$
estimates $\varepsilon=\sqrt{F_0/df}$. The difference from complete data is
entirely in \emph{which} $V$ enters $U$, and that is where the design leaks in.

\section{The target moves with the design}

Write $\lambda$ as a minimized quadratic form. From \eqref{eq:U},
$\tau'U\tau=\tau'V\tau-\tau'V\dot\sigma(\dot\sigma'V\dot\sigma)^{-1}\dot\sigma'V\tau$,
which is the $V$-orthogonal projection residual:
\begin{equation}\label{eq:lammin}
  \lambda = \min_{b\in\mathbb{R}^t}\,(\tau-\dot\sigma b)'\,V\,(\tau-\dot\sigma b).
\end{equation}
The structural error $\tau$ is a property of the population and the model, the
same vector whatever is missing. Only $V$ carries the missingness. Compare it to
the complete-data curvature $V_{\mathrm c}=\tfrac12 D'(\Sigma^{-1}\otimes\Sigma^{-1})D$
(bordered by the mean block).

\begin{lemma}[Information monotonicity]\label{lem:mono}
Each pattern's per-case information is the Fisher information of a coarsening
(marginalization) of the complete vector, so $\mathcal I_{o_i}\preceq V_{\mathrm c}$
in the Loewner order; averaging preserves it, $V=\plim\tfrac1N\sum_i\mathcal
I_{o_i}\preceq V_{\mathrm c}$. Under identification $V\succ0$.
\end{lemma}

\begin{theorem}[Missingness attenuates RMSEA]\label{thm:atten}
Fix the structural misspecification $\tau$. The FIML noncentrality is dominated
by its complete-data value,
\[
  \lambda_{\mathrm{FIML}} = \min_b(\tau-\dot\sigma b)'V(\tau-\dot\sigma b)
  \;\le\;
  \min_b(\tau-\dot\sigma b)'V_{\mathrm c}(\tau-\dot\sigma b)
  = \lambda_{\mathrm{complete}},
\]
hence the population RMSEA satisfies
$\varepsilon_{\mathrm{FIML}}\le\varepsilon_{\mathrm{complete}}$, with the gap
widening as the missing fraction grows.
\end{theorem}

\begin{proof}
By Lemma~\ref{lem:mono}, $(\tau-\dot\sigma b)'V(\tau-\dot\sigma b)\le
(\tau-\dot\sigma b)'V_{\mathrm c}(\tau-\dot\sigma b)$ for every $b$; the minimum
of a pointwise-smaller function is no larger. Divide by $n\,df$ and take roots.
\end{proof}

\begin{corollary}[FIML RMSEA is not design-invariant]\label{cor:noninvariant}
Two studies with identical population $(\mu_0,\Sigma_0)$ and identical model, but
different missingness, report different population RMSEAs; a misfit that reads
$0.08$ complete will read smaller under attrition, not because the model fits
better but because the missing data carry less information to expose the misfit.
RMSEA cutoffs calibrated on complete data therefore drift conservative under
missingness. A design-invariant rescaling would divide $\lambda$ by the
information fraction $\tr(V)/\tr(V_{\mathrm c})$ (a fraction-of-missing-information
correction), recovering $\varepsilon_{\mathrm{complete}}$; this is not what
lavaan or magmaan report, and reporting it would require the complete-data
curvature, which the incomplete sample does not deliver directly.
\end{corollary}

The attenuation is small under light, conditionally-benign missingness and grows
with the missing fraction; the empirical companion is the missing-data fit-index
literature (Davey \cite{davey2005}).

\section{Robust FIML RMSEA}

Drop normality, or admit MAR/observed-information effects: the identity $J=V$
breaks and $\Gmis=V^{-1}JV^{-1}\ne V^{-1}$. The expansion \eqref{eq:Tquad}
stands, but $U\Gmis$ is no longer idempotent.

\begin{lemma}[General FIML limit]\label{lem:mixture}
Under the drift with finite fourth moments,
$T\dto\sum_{j=1}^{df}\alpha_j\chi^2_{1,j}(\delta_j)$, the $\alpha_j$ the nonzero
eigenvalues of $U\Gmis$ and $\sum_j\alpha_j\delta_j=\tau'U\tau$. Under $J=V$ all
$\alpha_j=1$ and the mixture collapses to Theorem~\ref{thm:ncp}.
\end{lemma}

The exact-mean identity $\E[w'Uw]=\tr(U\Gmis)+\tau'U\tau$ again separates
sampling noise from misspecification additively:
\begin{equation}\label{eq:meandecomp}
  \E[T]\to \underbrace{\tr(U\Gmis)}_{=\,c\,df}+\underbrace{\tau'U\tau}_{=\,NF_0},
  \qquad
  c=\frac{\tr(U\Gmis)}{df}=\frac{\tr(UV^{-1}JV^{-1})}{df}.
\end{equation}
Here $c$ is the FIML Satorra--Bentler/Yuan--Bentler (MLR) scaling: $c=1$ under
$J=V$, and $c\ne1$ under non-normal or MAR-distorted scores. The
misspecification term $\tau'U\tau=NF_0$ is built from $U$ (hence $V$) alone and
does not see $J$, so it is the \emph{same} quadratic form as in
Theorem~\ref{thm:ncp}. The robust estimator follows by subtracting the inflated
mean:
\begin{equation}\label{eq:robust}
  \rmsea_{\mathrm{rob}}
    = \sqrt{\frac{\max(T-\hat c\,df,\,0)}{n\,df}}
    = \sqrt{\hat c}\;\rmsea_{\mathrm{sc}},
  \qquad
  \hat c=\frac{\tr(U\widehat\Gmis)}{df},
\end{equation}
with $\widehat\Gmis=\hat V^{-1}\hat J\hat V^{-1}$ from the sample casewise scores
and curvature (Yuan--Bentler \cite{yuanbentler2000}; the robust-RMSEA
construction of Brosseau-Liard--Savalei \cite{brosseauliardsavalei2012} and
Savalei \cite{savalei2018}). As in the complete-data note,
$\rmsea_{\mathrm{rob}}=\sqrt{\hat c}\,\rmsea_{\mathrm{sc}}$, and the scaled
version (which divides the noncentrality by $\hat c$ a second time) over-corrects;
\eqref{eq:robust} is the one that targets $\varepsilon$.

\begin{remark}[Observed vs.\ expected information: a FIML-only wrinkle]
With complete data and a correct normal model, the curvature $V$ is unambiguous.
Under missing data the \emph{observed} information (the realized casewise
Hessians) and the \emph{expected} information (their model expectation) differ in
finite samples even when $J=V$ asymptotically, and they diverge further under
non-normality (Savalei \cite{savalei2010}). The choice enters $\hat V$ in
\eqref{eq:robust}, hence $\hat c$, hence the robust RMSEA and its interval. lavaan
matches its MLR/robust output only under the observed-information,
structured-H1 default; the FMG spectrum convention reads the eigenvalues of
$U\Gmis$ off the \emph{unstructured} saturated H1. The point estimate is
robust to the choice through the mean identity \eqref{eq:meandecomp}; the
interval is not, and should declare its convention.
\end{remark}

\begin{remark}[Spectrum vs.\ mean: where the interval comes from]
The robust \emph{point} estimate rests only on \eqref{eq:meandecomp} and is
correct to leading order regardless of the $\alpha_j$. The interval inverts a
noncentral $\chi^2_{df}$ CDF, approximating the true mixture of
Lemma~\ref{lem:mixture} by a single noncentral $\chi^2$; the
mean-and-variance (Satterthwaite) refinement replaces $df$ by
$df^\ast=\tr(U\Gmis)^2/\tr\{(U\Gmis)^2\}$. The FMG/pEBA route used in the
companion experiments instead keeps the full eigenvalue spectrum of $U\Gmis$ and
forms the exact mixture reference, which is the right object when a few
$\alpha_j$ dominate and the single-$\chi^2$ approximation is poor. magmaan
should expose $\hat c$ (point), $df^\ast$ (Satterthwaite interval), and the
$U\Gmis$ spectrum (pEBA interval) from one FIML information pass.
\end{remark}

\begin{remark}[Truncation bites harder under FIML]
Because $\lambda_{\mathrm{FIML}}\le\lambda_{\mathrm{complete}}$
(Theorem~\ref{thm:atten}), the true noncentrality is smaller under missingness,
so $T-df$ (or $T-\hat c\,df$) lands below zero more often and $\rmsea$ piles up at
the $0$ floor. The negative-$\widehat{\mathrm{ncp}}$ events seen in the
experiment-23/24 cells are this finite-sample truncation artifact, not evidence
of inconsistency; they should shrink with $N$ and with the observed information,
and the interval (Corollary above), not the truncated point estimate, is the
honest summary.
\end{remark}

\begin{remark}[Two-stage and the completion trap]
The two-stage (ML2S) statistic first fills the saturated $\hat\eta_{\mathrm{sat}}$
then fits the structure to it as if complete, so its naive LRT uses the
\emph{complete-data} curvature $V_{\mathrm c}$ in place of $V$, inflating $T$ by
roughly the information ratio (a $\approx1.8\times$ over-rejection in the
experiment-24 MCAR-normal cell). Feeding that naive $T$ into the RMSEA formula
mislabels the inflation as misfit. The Satorra--Bentler correction with the
\emph{correct} two-stage sandwich (the Unstructured-H1 weight) restores the
$\chi^2_{df}$ mean and reproduces lavaan \texttt{robust.two.stage}; the robust
RMSEA \eqref{eq:robust} then transfers unchanged, with $\hat V$ the two-stage
curvature.
\end{remark}

\begin{remark}[Precondition: the target must exist]
All of the above roots $F_0$ as the population FIML discrepancy. Under MAR this
is the right target only when the incomplete-on-observed conditional means are
linear and conditional covariances homoscedastic (the FIML consistency note); if
those fail, $\hat\eta\not\pto\eta_0$, $F_0$ does not measure structural misfit,
and a robust correction repairs the reference law of a statistic that is no
longer centered on the model. No scaling fixes a moved target; that is an
estimator-bias problem, upstream of RMSEA.
\end{remark}

\begin{thebibliography}{9}

\bibitem{arminger1990}
Arminger, G., \& Sobel, M. E. (1990).
Pseudo-maximum likelihood estimation of mean and covariance structures with
missing data.
\emph{Journal of the American Statistical Association}, 85(409), 195--203.
\href{https://doi.org/10.1080/01621459.1990.10475324}{doi:10.1080/01621459.1990.10475324}.

\bibitem{brosseauliardsavalei2012}
Brosseau-Liard, P. E., Savalei, V., \& Li, L. (2012).
An investigation of the sample performance of two nonnormality corrections for
RMSEA.
\emph{Multivariate Behavioral Research}, 47(6), 904--930.
\href{https://doi.org/10.1080/00273171.2012.715252}{doi:10.1080/00273171.2012.715252}.

\bibitem{brownecudeck1992}
Browne, M. W., \& Cudeck, R. (1992).
Alternative ways of assessing model fit.
\emph{Sociological Methods \& Research}, 21(2), 230--258.
\href{https://doi.org/10.1177/0049124192021002005}{doi:10.1177/0049124192021002005}.

\bibitem{davey2005}
Davey, A. (2005).
Issues in evaluating model fit with missing data.
\emph{Structural Equation Modeling}, 12(4), 578--597.
\href{https://doi.org/10.1207/s15328007sem1204_4}{doi:10.1207/s15328007sem1204\_4}.

\bibitem{savalei2010}
Savalei, V. (2010).
Expected versus observed information in SEM with incomplete normal and nonnormal
data.
\emph{Psychological Methods}, 15(4), 352--367.
\href{https://doi.org/10.1037/a0020143}{doi:10.1037/a0020143}.

\bibitem{savalei2018}
Savalei, V. (2018).
On the computation of the RMSEA and CFI from the mean-and-variance corrected test
statistic with nonnormal data in SEM.
\emph{Multivariate Behavioral Research}, 53(3), 419--429.
\href{https://doi.org/10.1080/00273171.2018.1455142}{doi:10.1080/00273171.2018.1455142}.

\bibitem{savaleibentler2009}
Savalei, V., \& Bentler, P. M. (2009).
A two-stage approach to missing data: Theory and application to auxiliary
variables.
\emph{Structural Equation Modeling}, 16(3), 477--497.
\href{https://doi.org/10.1080/10705510903008238}{doi:10.1080/10705510903008238}.

\bibitem{steigershapirobrowne1985}
Steiger, J. H., Shapiro, A., \& Browne, M. W. (1985).
On the multivariate asymptotic distribution of sequential chi-square statistics.
\emph{Psychometrika}, 50(3), 253--263.
\href{https://doi.org/10.1007/BF02294104}{doi:10.1007/BF02294104}.

\bibitem{yuanbentler2000}
Yuan, K.-H., \& Bentler, P. M. (2000).
Three likelihood-based methods for mean and covariance structure analysis with
nonnormal missing data.
\emph{Sociological Methodology}, 30(1), 165--200.
\href{https://doi.org/10.1111/0081-1750.00078}{doi:10.1111/0081-1750.00078}.

\end{thebibliography}

\end{document}
